\chapter{Range Manipulation}

\section{Intervals}
TODO: explain mathematical definitions of intervals [a, b], (a, b) and in between
explain how to do operations on them, and what they meant intuitively...

Explain how intervals can be applied in functions to give a more complicated inter al

Explain the limitations of intervals in terms of mathematical expression; Intervals can't explain distribution of values within it.
It only explains the bounds.

\section{Useful Manipulations}
\subsection{Linear Mapping}
Given some value $x \in [a, b]$, we wish to map all values linearly to another range such that $x' \in [i, j]$,
we can follow this procedure(assuming that $a, b, i, j \in \mathbb{R}$):

\begin{enumerate}
    \item Shift $x$ from $[a, b]$ to $[0, (b - a)]$ by subtracting by $a$.
    \item Squeeze the value to range $[0, 1]$ by dividing by $(b - a)$.
    \item Scale the range $[0, 1]$ to $[0, (j - i)]$ by multiplying by $(j - i)$.
    \item Finally, offset the range $[0, (j - i)]$ to the desired range $[i, j]$ by adding $i$.
\end{enumerate}

In mathematical terms:
\begin{center}
    $m(x, a, b, i, j) = \frac{(x - a)(j - i)}{(b - a)} + i$
\end{center}

Very useful for mapping a value from one range to another!

This also works for ranges that are reversed(i.e. not following the $[min, max]$ format)! So for instance, let's map
from $[20, 30]$ to $[1, 0]$:
\begin{center}
    \begin{flalign}
        m(x, 20, 30, 1, 0) & = \frac{(x - 20)(0 - 1)}{(30 - 20)} + 1 \\
                           & = \frac{20 - x}{10} + 1 \\
                           & = \frac{20 - x + 10}{10} \\
                           & = \frac{30 - x}{10}
    \end{flalign}
\end{center}

We can check for it's validness by putting this in a graphing calculator like \textbf{Desmos}:
\begin{center}
    \includegraphics[width=25em, height=13em]{linear_mapping_reversed_example}
\end{center}

You can see $20$ points to $1$, and $30$ points to $0$ with a smooth linear line in between!

In fact, linear mapping formula used here is actually solving for a line formula,
where two points are $(20, 1)$ and $(30, 0)$ :D



\sgap
Here are some pre-derived formulas for common range mappings:

\tgap
From $x \in [a, b]$ to $x' \in [0, 1]$ (Also called the normal range)
\begin{center}
    $x' = \frac{x - a}{b - a}$
\end{center}

\sgap
From $x \in [a, b]$ to $x' \in [-1, 1]$ (Also called the signed normal range)
\begin{center}
    $x' = \frac{2(x - a)}{b - a} - 1$
\end{center}

\subsection{Value Clamping}
Sometimes we want to restrict some value $x \in \mathbb{R}$ within some max or a min boundary, we can utilize exactly
the $\max$ and $\min$ functions for this exact purpose! They are defined as follows:

\sgap
\begin{samepage}
    Here is the definition of the $\max$ function:
    \begin{center}
        \[
            \max(x, y) =
            \begin{cases}
                x & \text{if } x \geq y \\
                y & \text{if } x < y
            \end{cases}
        \]
    \end{center}
\end{samepage}

It gives you the max between two numbers $x, y \in \mathbb{R}$, pretty straight forward!

We can restrict values to not get below a certain value by utilizing this function, for instance
$\max(x, 20)$ forces the result to be greater of equal to $20$, in other words $\max(x, 20) \geq 20$. If $x$
goes below $20$, the $\max$ function always just picks $20$, otherwise it picks $x$ since it's greater.

\sgap

\begin{samepage}
    And this is the definition of $\min$ function:
    \begin{center}
        \[
            \min(x, y) =
            \begin{cases}
                x & \text{if } x \leq y \\
                y & \text{if } x > y
            \end{cases}
        \]
    \end{center}
\end{samepage}

Similar to $\max$ but reversed. It gives the minimum of two values $x, y \in \mathbb{R}$.

We can restrict values to not get above a certain value, for example:
$\min(x, 100)$, the result will never get above $100$, since if $x$ is greater than $100$, that would mean
$\min$ will just pick $100$ since it's smaller. Otherwise it'll pick $x$.


\sgap
Finally, we can utilize both to restrict a value $x \in \mathbb{R}$ within a certain range $[a, b]$ by using both one after the other.
This is often called \textbf{clamping}:
\begin{center}
    $clamp(x, a, b) = \min(\max(x, a), b)$
\end{center}

We first restrict $x$ to not get below $a$, then we restrict that result to not get above $b$ :D


\subsection{Value Stepping}
TODO


\subsection{Looping a Range}
If you have a value $x \in \mathbb{R}$, and you wish to do something to it, so it always stay in a range of $[a, b]$,
then you can follow the following procedure:

TODO





